% Part: lambda-calculus
% Chapter: introduction
% Section: conversion

\documentclass[../../../include/open-logic-section]{subfiles}

\begin{document}

\olfileid{lam}{int}{con}
\olsection{సమానరూప మార్పు మరియు సంక్షేపణం}

లాంబ్డా కలనంలో మూడు రకాల సమానరూప మార్పులు లేదా సంక్షేపణాలు ఉన్నాయి. తరువాత చూడబోతున్నట్లుగా, ``రూపాంతరం'' ఒకే దిశలోనా లేక రెండు దిశల్లోనా జరుగుతుందనేదే సమానరూప మార్పుకూ సంక్షేపణానికీ తేడా.

$\alpha$-సమానరూప మార్పు చరరాశుల పేరుమార్పుకు సంబంధించినది. $\lambda x. x$, $\lambda y.y$ ఒకే ప్రమేయాన్ని (తాదాత్మ్య ప్రమేయాన్ని) సూచిస్తున్నాయని భావించడం సహజం. ఈ ఆలోచనను కింది విధంగా ఔపచారికం చేస్తామని మూలం చెబుతుంది:
\sourcecorrection{OLTELAMSYNCONV-001}{మూల విభాగం ఈ పరిచయ వాక్యం తరువాతే ముగుస్తుంది; వాగ్దానం చేసిన $\alpha$-మార్పు నియమమూ, మిగిలిన రెండు మార్పుల వివరణా అందులో లేవు. మూలంలో లేని నియమాలను ఇక్కడ కల్పించి చేర్చలేదు.}

\end{document}
